\section{Exercise 6: Convolution}
\subsection{Task and method}
Implement \texttt{[yn, yorigin] = convolution(xn, xorigin, hn, horigin)}
to calculate
\begin{equation}
  y(n)=(x*h)(n)=\sum_{k=-\infty}^{\infty}x(k)h(n-k).
\end{equation}
For an LTI system with impulse response $h(n)$ and input $x(n)$,
this convolution gives the output. Only finitely many terms contribute
for the stored finite sequences, which are zero outside their intervals.

Let the input lengths be $N$ and $M$. Initialize an output vector with
$N+M-1$ zeros. For each pair of input vector positions $i$ and $j$,
accumulate its contribution at output position $i+j-1$:
\begin{lstlisting}[numbers=none]
for i = 1:N
    for j = 1:M
        yn(i+j-1) = yn(i+j-1) + xn(i)*hn(j);
    end
end
\end{lstlisting}
The product contributes at the sum of the input sample times,
\begin{equation}
  n=(i-\texttt{xorigin})+(j-\texttt{horigin}).
\end{equation}
At output vector position $r=i+j-1$, the time is also
$n=r-\texttt{yorigin}$. Equating these expressions gives
\begin{equation}
  \texttt{yorigin}=\texttt{xorigin}+\texttt{horigin}-1.
\end{equation}
Thus the output interval runs from $2-\texttt{xorigin}-\texttt{horigin}$
to $N+M-\texttt{xorigin}-\texttt{horigin}$. The implementation also
plots $x(n)$, $h(n)$, and $y(n)$ in a single figure.

\subsection{Source code}
\lstinputlisting[caption={Scilab convolution function and example call},label={lst:convolution}]{../6.sci}

\subsection{Example and expected result}
Choose the signal from Exercises 1--3 and a two-sample impulse response:
\begin{lstlisting}[numbers=none]
[yn, yorigin] = convolution([1, -2, 3, 6], 3, [1, 1], 1);
\end{lstlisting}
Here $x(n)$ is stored on $-2\leq n\leq1$ and
$h(n)=\delta(n)+\delta(n-1)$. Substituting this impulse response
into the convolution sum gives a direct check:
\begin{equation}
  y(n)=x(n)+x(n-1).
\end{equation}
\begin{table}[H]
  \centering
  \caption{Manual calculation of the convolution example}
  \label{tab:convolution}
  \begin{tabular}{lrrrrr}
    \toprule
    Time $n$ & $-2$ & $-1$ & 0 & 1 & 2\\
    \midrule
    $x(n)$ & 1 & $-2$ & 3 & 6 & 0\\
    $x(n-1)$ & 0 & 1 & $-2$ & 3 & 6\\
    $y(n)$ & 1 & $-1$ & 1 & 9 & 6\\
    \bottomrule
  \end{tabular}
\end{table}
The calculated result is
\[
  \texttt{yn}=[1,-1,1,9,6],\qquad
  \texttt{yorigin}=3+1-1=3.
\]
The output has $4+2-1=5$ stored samples. Its third element corresponds
to $n=0$.

\subsection{Execution evidence}
\resultfigure{images/exercise6.png}
  {Input, impulse response, and convolution output for Exercise 6}
  {fig:convolution}
The captured Scilab plot agrees with \Cref{tab:convolution}: the output
samples are $1,-1,1,9,6$ at $n=-2,-1,0,1,2$, respectively.

\subsection{Discussion}
Each output sample is a sum of products at the appropriate time offsets.
For this impulse response, the system adds the current input sample and
the preceding input sample. In particular, $y(1)=6+3=9$ and
$y(2)=0+6=6$. The output therefore extends one sample beyond the
stored input interval. The zero padding used to describe the finite
sequences must be respected when calculating these boundary samples.
